			One of the main objectives of empirical economics is identifying causal effects. Identification strategies such as Regression Discontinuity (RD), Instrumental Variables (IV), Difference-in-Differences (DiD) and event studies help us achieve this goal. To isolate causal effects, these strategies only use a portion of the variation in the data. They exploit the exogenous part of the variation in the treatment or decrease the sample size by only considering observations for which the as-if random assignment assumption is credible. Limiting the variation used can decrease precision and thus statistical power---the probability of rejecting the null hypothesis when it is false, or put simply, the probability of obtaining a statistically significant estimate. There is, therefore, a tension between reducing confounding and statistical power. Yet, when statistical power is low, not only is the estimator imprecise but statistically significant estimates exaggerate the true effect size \citep{gelmanType2000, ioannidis_why_2008, gelman_beyond_2014}. For statistically significant estimates, the tension is therefore between reducing confounding and exaggerating the true effect size. While causal identification strategies alleviate confounding bias, this paper argues that they can introduce another form of bias---exaggeration---by limiting variation used for identification. 
			
			% what I do in this paper
			This paper aims to demonstrate the existence and practical importance of this confounding-exaggeration trade-off. To do so, I analyze published results from the literature, derive the trade-off theoretically in a simple setting, and explore its main drivers through Monte Carlo simulations. I then discuss practical solutions to navigate this trade-off and introduce a tool to directly examine the overarching mechanism behind it: the identifying variation and which observations actually contribute to estimation. While causal identification strategies are essential for describing causal relationships, the present paper emphasizes that a perfectly convincing identification does not guarantee an absence of bias. Paradoxically, strengthening identification can actually pull estimates away from the true effect by introducing another type of bias.
			
			Exaggeration of significant estimates arises in imprecise and underpowered studies because such estimates make up a selected sub-sample of all estimates. Those statistically significant at the 5\% level are at least 1.96 standard errors away from zero and if these standard errors are large, significant estimates are located in the tails of the distribution of all possible estimates. %The average of these statistically significant estimates differs from the true effect, located at the center of the distribution if the estimator is unbiased. 
			They exaggerate the true effect and the less precise the estimator, the larger exaggeration is.  \Cref{fig:graph_exag} illustrates the inflation of significant estimates caused by imprecision. When power is low, obtaining a statistically significant estimate from an unbiased estimator guarantees that it will be far from the true effect. An estimator $\hat{\beta}$ of the true effect $\beta$ might be unbiased in the traditional sense of $\mathbb{E}[\hat{\beta}] = \beta$ but conditionally biased in the sense that $\mathbb{E}[\hat{\beta} | \text{ Significant}] \neq \beta$. 
			
			   \begin{figure}[!h]
				\begin{center}
					\caption{Significance and distribution of two unbiased estimators with different variances}
					\label{fig:graph_exag}
					\includegraphics[width=0.8\linewidth]{images/graph_intuition_precision.pdf}
					\caption*{\footnotesize \textit{Notes}: 100,000 draws from two normal distributions $\mathcal{N}(1, 0.05)$ and $\mathcal{N}(1, 0.5)$. Significance level: 5\%}
				\end{center}
				\vspace{-1cm}
			 \end{figure}
			
			Exaggeration only arises under two conditions: 1) a publication bias favors a certain type of results and 2) statistical power is low. Both features are observed in the economics literature. A large body of literature underlines selection on statistical significance in economics \citep[for instance]{rosenthal_file_1979, brodeur_star_2016, andrews_identification_2019, abadie_statistical_2020, brodeur_methods_2020}. Additional studies have highlighted its frequent and substantial lack of statistical power and resulting exaggeration \citep{ioannidis_power_2017, ferraro_featureis_2020}. Even in experimental economics, with a high level of control and an arguable absence of confounders, studies from top economics journals failed to replicate, the original estimates being on average inflated by a factor of at least 1.5 \citep{camerer_evaluating_2016}. In the non-experimental economics literature, where statistical power is rarely a central consideration under current practices, several meta-analyses provide evidence of consequential exaggeration. \cite{ioannidis_power_2017} estimates the median statistical power in a wide range of areas of economics to only 18\%. Despite %the widespread use of convincing causal identification strategies and 
			usually large sample sizes, they show that nearly 80\% of estimates are likely exaggerated by a factor of two. As further illustrated in Section \ref{lit_review}, the magnitude of exaggeration can be considerable and in some situations could be on par with that of a bias caused by confounders.
			%especially in settings where bias caused by confounders may not be as large. For example, Weidmann and Miratrix (2021) documents and absence of evidence of selection bias due to unobservables in the evaluations of school programs they investigate.
			Taking exaggeration into account and understanding its drivers is therefore crucial.
			
			Accurate point estimates are instrumental as they often inform policy decisions through Cost-Benefit Analyses (CBA). For instance, environmental economics estimates enter the computation of the Social Cost of Carbon or routinely help policy makers decide on the implementation of regulations. Yet, the underlying effects can be relatively small and thus difficult to capture, making the studies subject to exaggeration. Estimates of the impact of environmental regulations on job losses constitute a good example \citep{gray_environmental_2023}. \cite{walker_environmental_2011} documents the impact of the Clean Air Act amendments of 1990 on employment and finds an effect of -14.2\% (s.e. 4.3). For similar policies, other studies find smaller effects, of the order of magnitude of -3\% \citep{greenstone_impacts_2002, gray_environmental_2023}. If the true effect was in fact of this magnitude, the statistical power of the design in \cite{walker_environmental_2011} would be 11\% and significant estimates would exaggerate the true effect by a factor of 3.5 on average. In this example, bias caused by exaggeration would be substantial and could have detrimental policy implications.\footnote{Of course, in the particular setting of \cite{walker_environmental_2011}, the true effect might be larger than 3\%. I am not claiming that the study is flawed but instead that its level of precision would produce inflated significant estimates if the true effect was of the order of magnitude of 3\%.}
			
			%mechanisms
			All causal identification strategies discard variation in order to identify causal effects but the confounding-exaggeration trade-off is mediated through a distinctive channel for each of them. RD designs discard part of the variation by only considering observations within the bandwidth, decreasing the effective sample size and thus precision. An IV strategy only uses the subset of the variation in the treatment that is explained by the instrument. In studies leveraging exogenous shocks, the variation used to identify an effect sometimes only comes from a limited number of treated observations. Matching prunes units that cannot be matched and thus reduces the effective sample size. Adding controls or fixed effects can increase the variance of the estimator and exaggeration if they absorb more of the variation in the treatment than in the outcome variable. 
			Since causal identification strategies can be interpreted as ways of controlling for confounders, this last point actually ties all the strategy-specific arguments together.\footnote{Fixed Effects (FEs) based identification strategies such as DiD control for the invariant, unobserved, and arguably endogenous part of the variation in the outcome. An IV approach essentially partials out the variation in $x$ unexplained by the instruments. Fuzzy-RD and propensity score matching can be thought of as control function approaches, of the forcing variable and propensity score respectively. In addition, excluding observations that are outside the bandwidth or unmatched is equivalent to controlling for observation-level fixed effects for these observations.} When these identification strategies absorb more of the variation in the treatment than in the outcome, they increase the variance of the resulting estimator and can cause exaggeration. Considering a simple linear homoskedastic model gives the intuition for this trade-off between exaggeration and omitted variable bias (OVB) for control approaches.  Let $y_{i} = \alpha + \beta x_{i} + \delta w_{i} + u_{i}$, $\forall i \in \{1, .., n\}$,  with $x$ the variable of interest, $w$ a potentially unobserved variable correlated with $x$ and $u$, an error term. Under usual assumptions and using the Frisch-Waugh-Lovell theorem, we get that $ \mathbb{V}_{\hat{\beta}_{\textsc{ovb}}}$ and $ \mathbb{V}_{\hat{\beta}_{\textsc{ctrl}}}$, the variance of the estimators for $\beta$ when omitting $w$ (short regression) and controlling for it (long regression) are respectively:
			~
			\[
				\mathbb{V}_{\hat{\beta}_{\textsc{ovb}}} =
				 \dfrac{\sigma^{2}_{u_{\textsc{ovb}}}}{n \ \sigma_{x}^{2}} =
				 \dfrac{\sigma^{2}_{y^{\perp x}}}{n \ \sigma_{x}^{2}}
				 \qquad \text{and} \qquad
				 \mathbb{V}_{\hat{\beta}_{\textsc{ctrl}}} = 
				 \dfrac{\sigma^{2}_{u_{\textsc{ctrl}}}}{n \ \sigma_{x^{\perp w}}^{2}} =
				  \dfrac{\sigma^{2}_{y^{\perp x, w}}}{n \ \sigma_{x^{\perp w}}^{2}}
			\]
			
			where $\sigma^{2}_{u_{\textsc{ovb}}}$ and $\sigma^{2}_{u_{\textsc{ctrl}}}$ are the variances of the residuals in the regression of $y$ on $x$ and of $y$ on $x$ and $w$ respectively or equivalently the variances of the parts of $y$ that are orthogonal to $x$ and to $x$ and $w$ respectively ($\sigma^{2}_{y^{\perp x}}$ and $\sigma^{2}_{y^{\perp x, w}}$), $\sigma^{2}_{x}$ is the variance of $x$ and $\sigma^{2}_{x^{\perp w}}$ is the variance of the part of $x$ orthogonal to $w$. Thus,
			
			\[
				\mathbb{V}_{\hat{\beta}_{\textsc{ovb}}} <  \mathbb{V}_{\hat{\beta}_{\textsc{ctrl}}}
                   		\quad \Leftrightarrow \quad 
				\dfrac{\sigma^{2}_{y^{\perp x}}}{n \ \sigma_{x}^{2}} < \dfrac{\sigma^{2}_{y^{\perp x, w}}}{n \ \sigma_{x^{\perp w}}^{2}}
				\quad \Leftrightarrow \quad 
				\dfrac{\sigma^{2}_{x^{\perp w}}}{\sigma_{x}^{2}} <
				\dfrac{\sigma^{2}_{y^{\perp x, w}}}{ \sigma_{y^{\perp x}}^{2}}
			\]
			
			Controlling for $w$ increases the variance of the estimator if the fraction of the variance unexplained by $w$ is greater for $y^{\perp x}$ than for $x$. Put differently, if controlling absorbs more of the variation in $x$ than in the residual part of $y$ ($y^{\perp x}$), the variance of the estimator increases. Since exaggeration increases with estimator's variance, controlling for a confounder can increase exaggeration. Controls, including fixed effects, limit the identifying variation and may cause identification to rely on a small effective sample. Observations whose treatment status is largely explained by these controls contribute little, if at all, to identification. Under some circumstances discussed in Sections \ref{maths} and \ref{simulations}, controlling can even produce an exaggeration bias larger than the OVB that would result from an absence of controls. 
			%This has direct implications when using fixed effects: an exaggeration bias can arise from the use of fixed effects if they absorb more of the variation in $x$ than in $y^{\perp x}$. %FIND EXAMPLES !! ---which can often be the case in settings such as regression of agricultural yield on weather variables or ---
			
			% what we actually do
			In the remainder of the paper, I first document the magnitude of the trade-off. To do so, I build on existing literature reviews \citep{brodeur_methods_2020, young_consistency_2022, bagilet_accurate_2023, lalHow2024} to discuss evidence of exaggeration bias in a large set of causal studies mostly published in top journals. %I then focus on IV designs as they allow for a within-study comparison between causal and non-causal estimates, the Two-Stage Least-Squares (2SLS) and the corresponding  ``naive'' Ordinary Least Squares (OLS) respectively. %IVs often reduces precision, limits statistical power to the point that 2SLS estimates reject the OLS point estimate and are not able to provide strong evidence supporting that the OLS is biased \citep{youngConsistencyInferenceInstrumental2022}. 
			These analyses reveal heterogeneity across studies: while exaggeration might be limited in some analyses, it is likely substantial in many others. 
			%They would not be able to accurately retrieve effect size of the magnitude of the OLS estimate and significant results would overestimate it.
			%By reducing precision, IVs---and causal identification strategies in general---likely induce exaggeration and produce biased significant estimates: 
			For instance, in the set of papers reviewed in \cite{young_consistency_2022}, I show that 
			%less than 10\% of the designs would retrieve an effect size equal to that of the ``naive'' OLS---at the conventional 80\% power threshold. This proportion increases to 17\% for headline results. 
			half of the IV designs would exaggerate a true effect size of the magnitude of the ``naive'' OLS estimate by a factor larger than 3.2---and by a factor larger than 2.0 for headline designs. %Similarly, in the literature on the short term health effects of air pollution that I review in \cite{bagilet_accurate_2023}, only 2\% of the IV designs would retrieve an effect size equal to that of the ``naive'' OLS and the median exaggeration factor would be 4.5. 
			I then directly compare confounding and exaggeration biases for an example study. I show that in this instance the bias of the IV could be on par or even larger than that of the OLS and that exaggeration may explain this difference. 
			
			Next, I derive a formal proof of the existence of the trade-off for prevailing causal identification strategies. Specifically, I show that the bias caused by exaggeration can be larger than the one caused by confounders. I also analyze the drivers of exaggeration and demonstrate that it increases as the strength of the instruments decreases, the number of exogenous shocks decreases or when controlling for a confounder absorbs more of the variation in the treatment than in the outcome. 
			
			 Then, I show that this \textit{causal exaggeration} arises for realistic parameter values by further exploring its drivers in realistic settings in which there is no closed form formula for power and exaggeration available. Exaggeration being defined with respect to the true effect, a quantity which is never known in real world settings, I turn to simulations to be able to compute it. Monte-Carlo simulations also allow varying the value of the parameter of interest \textit{ceteris paribus}. To make these simulations concrete, I calibrate them to emulate existing studies from education, environmental, health, labor, and political economics. %All other simulation assumptions are chosen to facilitate the recovery of the effect of interest: I consider simple linear models with constant and homogenous treatment effects, \textit{i.i.d.} observations and homoskedastic errors. All the models are correctly specified and accurately represent the data generating process, except for the omitted variable. 
			I find that causal exaggeration can be substantial in various realistic settings where the variation remaining for identification is limited. 
									
			Finally, I discuss concrete avenues to address causal exaggeration when carrying out a non-experimental study.\footnote{In experimental studies, there are essentially no confoundings. A solution to increase power and reduce exaggeration is generally to increase sample size, reduce noise by improving measurement or improving balance or to focus on larger potential effects.} A series of tools can be used to evaluate the potential magnitude of confounding and exaggeration issues separately. Sensitivity analyses help with the former while power calculations help with the latter. %These power calculations can be computed before and after the analysis is carried out. 
			%For instance, the sensitivity analysis tools developed in \cite{cinelli_making_2020} enable to assess how strong confounders would have to be to change the estimate of the treatment effect beyond a given level we are interested in. 
			Considering the attention given to bias avoidance in the economics literature, I underline that making power central to non-experimental analyses, even after an effect has been found, would help limiting bias caused by exaggeration. Prospective power simulations help identify the design parameters affecting power and exaggeration by approximating the data generating process \citep{gelman_regression_2020, black_simulated_2022}. Retrospective power calculations allow evaluating whether a study would have had enough power to confidently estimate a range of smaller but credible effect sizes \citep{gelman_beyond_2014, stommesReliability2023}.
			%%%%%%%% Evaluating potential exaggeration and confounders is however complex as both rely on unobserved stuff
			Focusing more specifically on the trade-off and its drivers, I consider and develop tools to identify the variation actually used for identification when using causal identification strategies. %The \href{https://vincentbagilet.github.io/causal_exaggeration/}{companion website} describes in details how such analyses can be implemented. Finally, I briefly discuss potential solutions to mitigate this trade-off.
			I introduce a \verb?R? package, \href{https://vincentbagilet.github.io/ididvar/}{\verb?ididvar?}, providing a set of tools to easily identify the identifying variation in a regression and the set of observations actually contributing to the estimation of the effect of interest.
						
			% Contributions
			This paper contributes to four strands of the applied economics literature. First, the idea that causal identification estimators, while unbiased, may be imprecise is not new; it is part of the well-known bias–variance trade-off \citep{imbens_optimal_2012, deaton_understanding_2018, hernan_causal_2020, ravallion_should_2020a}. I approach this literature from a different angle; through the prism of statistical power and publication bias. Limited precision resulting from the use of causal identification strategies not only makes it difficult to draw clear conclusions about the exact magnitude of effects, but, I argue, can also inherently lead to inflated published estimates, creating another form of bias. In this sense, the bias–variance trade-off is in fact a bias–bias trade-off.
			
			Second, studies discussing the exaggeration of statistically significant estimates due to low power rarely investigate its determinants or consider specific causal identification strategies separately \citep{ioannidis_power_2017, schellEvaluatingMethodsEstimate2018, ferraro_featureis_2020, black_simulated_2022, stommesReliability2023, arel-bundock_quantitative_2024}. In a companion paper, I highlight tangible design parameters that can generate exaggeration across a wide range of research designs \citep{bagilet_accurate_2023}. Other common empirical practices in economics are also associated with larger standard errors and, consequently, greater risks of exaggeration. For instance, clustering, while essential for accounting for unspecified and potentially unobserved correlation structures in the error terms, tends to increase standard errors. \citet{vu_standard_2025} examines this in detail. In the present paper, I take a step back and propose an overarching mechanism inherent to causal identification strategies as a whole: although each strategy does so through different means, in essence they discard part of the variation, thereby increasing the risks of exaggeration.
			
			Third, this paper makes a methodological and practical contribution by facilitating the computation of concrete measures of each observation’s contribution to identification and, by extension, of the effective sample underlying the estimated effect. It extends and complements existing approaches \citep{angrist_mostly_2009, aronow_does_2016} and facilitates their implementation and interpretability in applied economics research. %One of the main metrics is an observation level identifying weight that corresponds to its leverage but adapted to the specificities of applied economics where we are typically interested in one parameter and not the whole vector of parameters, \textit{ie} after partialling out controls. Another one identify a set of observations that can be ``dropped'' without affecting the point estimate and standard error by more than a small amount. 
			The accompanying \verb?R? package enables users to compute, visualize, and analyze these measures in a straightforward and consistent way. It allows applied researchers to diagnose potential exaggeration risks \textit{ex-post}, without assumptions about true effect sizes. Beyond the scope of this paper, the package also provides a framework for identifying and communicating the set of observations contributing to identification and how this set changes across specifications, levels of fixed effects and set of control variables.
			
			Finally, this study contributes to the literature on replicability in economics \citep{camerer_evaluating_2016, ioannidis_power_2017, christensen_transparency_2018, kasy_forking_2021}. The trade-off presented in this paper suggests that the widespread reliance on convincing causal identification strategies in economics does not necessarily shield the field from potential replication threats. In addition, the tools discussed and introduced can be readily used in replication analyses or referee reports to transparently document statistical power and the sources of identifying variation and can help assess robustness across different specifications, subsamples, or sets of controls.
			
			%outline
			In the following section, I document evidence of causal exaggeration in the causal  economics literature. In Section \ref{maths}, I study the drivers of exaggeration and formally show in a simple setting that the use of causal identification strategies can exacerbate it. In Section  \ref{simulations}, I implement realistic Monte-Carlo simulations to illustrate the existence of the confounding-exaggeration trade-off. I discuss potential solutions to navigate this trade-off and introduce the \verb?ididvar? package in Section \ref{discussion} and conclude in Section \ref{conclusion}.
	
	
		